\documentclass[11pt,a4paper]{article}
\usepackage[utf8]{inputenc}
\usepackage[T1]{fontenc}
\usepackage[brazil]{babel}
\usepackage[a4paper,margin=2.3cm]{geometry}
\usepackage{enumitem}
\usepackage{hyperref}
\usepackage{xcolor}
\setlength{\parindent}{0pt}
\setlength{\parskip}{5pt}
\sloppy
\setlength{\emergencystretch}{3em}
\hypersetup{colorlinks=true,urlcolor=blue,linkcolor=black}
\newcommand{\cabecalho}[2]{\begin{center}{\small \textbf{Sprint 2 -- #1}}\\[1mm]\textit{Inventário de Ativos e Vulnerabilidades}\\Matheus Moreli Durão Silva -- 12621CBS217\\Cibersegurança -- Bacharelado -- Noturno\end{center}\hrule\vspace{4mm}}
\newcommand{\bloco}[1]{\par\vspace{2mm}\textbf{#1}\par}
\newcommand{\codigo}[1]{\texttt{#1}}

\begin{document}\cabecalho{S2\_19 -- Entrega do Trabalho 1}{}
\bloco{Conteúdo da entrega} Código particionado, JSON, especificação, critérios de aceitação, documentação dos cards e testes. O arquivo \codigo{app.py} é o ponto de entrada; os demais módulos explicam a divisão de responsabilidade.
\bloco{Como executar} Abrir um terminal na pasta do projeto e usar \codigo{python app.py}. Para verificar as regras sem interagir com o menu, usar \codigo{python -m unittest -v test\_sprint2.py}. Não há dependência de biblioteca de terceiros.
\bloco{Auditoria} A especificação S2\_07 define escopo e modelo; S2\_04 e S2\_06 ligam requisitos a critérios; \codigo{test\_sprint2.py} oferece evidência executável. O arquivo JSON é legível e não contém campos de credenciais.
\bloco{Pendências externas antes de finalizar no GitHub} Conferir links e permissões do repositório; criar tag da versão final; usar branches/PRs; obter revisão real de colega para S2\_17. Esses itens dependem de ações no GitHub, não podem ser marcados como concluídos apenas pelo código local.
\bloco{Checklist} Aceitação e testes cobertos; instrução de execução disponível; sem segredos; pendências de links/permissões/tag explicitadas para conclusão no repositório.
\end{document}
